% ============================================================
%  Chapitre 2 — Étude Préliminaire
% ============================================================
\chapter{Étude Préliminaire}
\thispagestyle{fancy}

% ─────────────────────────────────────────────────────────────
\section{Introduction}
% ─────────────────────────────────────────────────────────────

Ce chapitre constitue la phase d'analyse préalable au développement de \farmAI{} v3.0.
Il commence par l'identification précise des besoins fonctionnels et non-fonctionnels
du système, puis dresse une revue des travaux académiques récents dans les domaines
de la vision par ordinateur agricole, de l'IoT et des LLM. L'architecture globale
et les choix technologiques sont ensuite présentés et justifiés, avant de détailler
la planification du projet par sprints Scrum.

% ─────────────────────────────────────────────────────────────
\section{Identification des besoins}
% ─────────────────────────────────────────────────────────────

\subsection{Les besoins fonctionnels}

Les besoins fonctionnels décrivent ce que le système \textbf{doit faire} — les
services rendus aux utilisateurs et aux acteurs externes. Ils sont issus de
l'analyse des cas d'utilisation identifiés au chapitre 1.

\begin{table}[H]\centering
\caption{Besoins fonctionnels — Smart Farm AI v3.0}
\label{tab:bf_all}
\small
\begin{tabular}{cp{4.5cm}p{6.5cm}l}
\toprule
\textbf{ID} & \textbf{Module} & \textbf{Description} & \textbf{Priorité} \\
\midrule
BF01 & Authentification & Inscription/connexion par email + mdp hashé (bcrypt) & Must \\
BF02 & Authentification & Double facteur OTP email (6 chiffres, 10 min) & Must \\
BF03 & Authentification & JWT access token (30j) + refresh token (7j) & Must \\
BF04 & Gestion Fermes   & CRUD complet fermes avec coordonnées GPS & Must \\
BF05 & Gestion Fermes   & Tableau de bord KPIs temps réel multi-fermes & Must \\
BF06 & IoT Télémétrie   & Réception métriques ESP32 via HTTP POST & Must \\
BF07 & IoT Télémétrie   & Visualisation courbes historiques (1h/24h/7j) & Must \\
BF08 & IoT Alertes      & Détection anomalies IsolationForest en temps réel & Must \\
BF09 & IoT Alertes      & Notifications WebSocket (SWARM\_RISK, THERMAL) & Must \\
BF10 & Vision IA        & Inférence YOLO sur image uploadée ou caméra live & Must \\
BF11 & Vision IA        & Choix parmi 8 modèles et 40 classes agricoles & Must \\
BF12 & Vision IA        & Historique des scans avec résultats & Should \\
BF13 & Apiculture       & CRUD ruchers, ruches avec QR Code & Must \\
BF14 & Apiculture       & Enregistrement visites et observations & Must \\
BF15 & Apiculture       & Suivi récoltes, stocks, dépenses & Should \\
BF16 & Apiculture       & Planning de traitement et rappels & Should \\
BF17 & Apiculture       & Télémétrie NODE\_B (poids, temp. ruche) & Must \\
BF18 & Assistant IA     & Conversation en Darija (texte + voix) & Must \\
BF19 & Assistant IA     & RAG sur corpus UTAP/AVFA & Must \\
BF20 & Assistant IA     & Analyse d'image multimodale (LLaVA) & Should \\
BF21 & Ouvriers         & CRUD ouvriers et assignation de tâches & Should \\
BF22 & PWA Worker       & Interface mobile offline (Dexie.js) & Must \\
BF23 & Cartographie     & Carte SIG Leaflet avec fermes et vétérinaires & Should \\
BF24 & MLOps            & Dashboard métriques MLflow (Silhouette, anomalie) & Could \\
\bottomrule
\end{tabular}
\end{table}

\noindent Les priorités suivent la classification \textbf{MoSCoW} :
\begin{itemize}
  \item \textbf{Must} (Obligatoire) : fonctionnalités indispensables à la valeur
        minimale du produit.
  \item \textbf{Should} (Important) : fortement souhaitées mais non bloquantes.
  \item \textbf{Could} (Optionnel) : ajoutées si le temps le permet.
\end{itemize}

\subsection{Les besoins non fonctionnels}

Les besoins non fonctionnels définissent les qualités que le système doit exhiber
— performance, sécurité, disponibilité, maintenabilité — indépendamment des
fonctionnalités métier.

\begin{table}[H]\centering
\caption{Besoins non fonctionnels — Smart Farm AI v3.0}
\label{tab:bnf}
\begin{tabular}{llll}
\toprule
\textbf{Catégorie} & \textbf{Besoin} & \textbf{Cible} & \textbf{Vérifié par} \\
\midrule
Performance   & Latence inférence YOLO      & $<$ 200 ms (p95)   & Benchmark pytest \\
Performance   & Temps réponse API (p95)     & $<$ 500 ms         & locust load test \\
Disponibilité & Uptime service backend      & $\geq$ 99\%        & Oracle monitoring \\
Sécurité      & Hashage mots de passe       & bcrypt cost=12     & Revue de code \\
Sécurité      & Expiration tokens           & Access 30j / OTP 10min & Tests unitaires \\
Sécurité      & Isolation multi-tenant      & farm\_id dans chaque requête & Tests pytest \\
Offline       & PWA Worker hors-ligne       & IndexedDB Dexie.js & Test réseau coupé \\
Souveraineté  & LLM exécuté localement      & Ollama local (no cloud) & Revue archi \\
Scalabilité   & Multi-fermes par user       & Illimité (FK isolée) & Tests DB \\
Maintenabilité & Couverture tests backend   & $\geq$ 80\%        & pytest --cov \\
Déploiement   & Cloud gratuit               & Oracle Free ARM64  & Déploiement effectif \\
IoT           & Fréquence télémétrie        & 10 secondes (configurable) & Firmware \\
\bottomrule
\end{tabular}
\end{table}

% ─────────────────────────────────────────────────────────────
\section{Revue des travaux récents}
% ─────────────────────────────────────────────────────────────

Cette section dresse un état de l'art ciblé sur les trois axes technologiques de
\farmAI{} : la détection de pathologies agricoles par vision par ordinateur,
la surveillance IoT agricole, et les assistants conversationnels en langue locale.

\begin{table}[H]\centering
\caption{Tableau comparatif des travaux récents — Vision IA agricole}
\label{tab:art_vision}
\small
\begin{tabular}{p{3.5cm}p{2cm}p{2.5cm}p{3cm}p{2.5cm}}
\toprule
\textbf{Auteurs} & \textbf{Année} & \textbf{Modèle} & \textbf{Application} & \textbf{Résultat} \\
\midrule
Kamilaris \& Prenafeta-Boldú \cite{kamilaris2018deep} & 2018 & CNN (AlexNet/VGG) & Détection maladies feuilles & mAP 95\% \\
Liakos et al. \cite{liakos2018machine} & 2018 & Random Forest / SVM & Prédiction rendements & R²=0.87 \\
Ultralytics \cite{ultralytics2023yolov8} & 2023 & YOLOv8 & Détection multi-classe & mAP50=89.2\% \\
\textbf{\farmAI{} v3.0} & \textbf{2025} & \textbf{YOLOv11 (8 modèles)} & \textbf{Pathologies multi-espèces} & \textbf{mAP50=0.891} \\
\bottomrule
\end{tabular}
\end{table}

\begin{table}[H]\centering
\caption{Tableau comparatif des travaux récents — IoT et détection d'anomalies}
\label{tab:art_iot}
\small
\begin{tabular}{p{3.5cm}p{2cm}p{2.5cm}p{3cm}p{2.5cm}}
\toprule
\textbf{Auteurs} & \textbf{Année} & \textbf{Algorithme} & \textbf{Données} & \textbf{Résultat} \\
\midrule
Liu et al. \cite{liu2008isolation} & 2008 & IsolationForest & Séries temporelles & AUC=0.92 \\
Sculley et al. \cite{sculley2015hidden} & 2015 & Systèmes ML prod & Pipelines données & Dette tech ML \\
\textbf{\farmAI{} v3.0} & \textbf{2025} & \textbf{IsolationForest} & \textbf{IoT ESP32 (7 features)} & \textbf{Silhouette=0.73} \\
\bottomrule
\end{tabular}
\end{table}

\begin{table}[H]\centering
\caption{Tableau comparatif des travaux récents — LLM et NLP en langue arabe}
\label{tab:art_llm}
\small
\begin{tabular}{p{3.5cm}p{2cm}p{2.5cm}p{3cm}p{2.5cm}}
\toprule
\textbf{Modèle/Travail} & \textbf{Année} & \textbf{Langue} & \textbf{Tâche} & \textbf{Limite} \\
\midrule
AraBERT & 2020 & Arabe MSA & Classification & Pas de Darija \\
CAMeL-LM & 2022 & Arabe dialectal & NLP & Licence NC \\
Lewis et al. \cite{lewis2020rag} & 2020 & Anglais & RAG Knowledge & Pas multilingue \\
Ghezaiel \cite{labess7b} & 2024 & \textbf{Darija tunisien} & \textbf{Chat local} & Latence CPU \\
\textbf{\farmAI{} v3.0} & \textbf{2025} & \textbf{Darija + FR + AR} & \textbf{RAG agricole} & \textbf{15-25s CPU} \\
\bottomrule
\end{tabular}
\end{table}

\begin{tcolorbox}[infobox, title={Positionnement de \farmAI{} dans la littérature}]
\farmAI{} v3.0 se distingue par \textbf{trois contributions originales} non couverts
par les travaux existants :
\begin{enumerate}
  \item Combinaison simultanée de 8 modèles YOLO v11, IoT ESP32 et RAG Darija dans
        un même système intégré.
  \item Premier assistant agricole souverain en dialecte tunisien, exécuté localement
        sans dépendance cloud externe.
  \item Application des pratiques MLOps industrielles (DVC + MLflow) au contexte
        agritech tunisien sur matériel accessible (CPU, Oracle Free Tier).
\end{enumerate}
\end{tcolorbox}

% ─────────────────────────────────────────────────────────────
\section{Architecture du projet}
% ─────────────────────────────────────────────────────────────

\farmAI{} v3.0 adopte une architecture en \textbf{5 couches} inspirée du modèle
N-tiers, garantissant la séparation des responsabilités et la testabilité indépendante
de chaque composant.

\begin{figure}[H]\centering
\begin{tikzpicture}[x=1cm, y=1cm,
  layer/.style={rectangle, rounded corners=5pt, draw=#1!70!black, fill=#1!10,
    minimum width=14cm, minimum height=1.3cm, align=center,
    font=\small\bfseries\color{#1!80!black}}]

\node[layer=farmblue]   at (0, 0)   {\textbf{Couche 1 — Clients} \quad
  Owner Dashboard (React SPA) \quad Worker PWA \quad Map Center \quad Assistant};

\node[layer=farmgreen]  at (0,-2)   {\textbf{Couche 2 — API Gateway} \quad
  FastAPI 50+ Endpoints \quad JWT Middleware \quad WebSocket \quad CORS};

\node[layer=farmorange] at (0,-4)   {\textbf{Couche 3 — Services Métier} \quad
  CV/YOLO Service \quad RAG/LLM Agent \quad Anomaly Service \quad Weather};

\node[layer=farmred]    at (0,-6)   {\textbf{Couche 4 — Données} \quad
  PostgreSQL \quad ChromaDB \quad MLflow Registry \quad DVC \quad Ollama};

\node[layer=farmgray]   at (0,-8)   {\textbf{Couche 5 — IoT} \quad
  NODE\_A (Irrigation ESP32) \quad NODE\_B (Rucher ESP32) \quad Wokwi Sim.};

\foreach \y/\lbl in {-0.65/HTTPS REST, -2.65/Appels Python, -4.65/SQLAlchemy, -6.65/HTTP POST}
  \draw[->, >=Stealth, thick, farmgray!70] (0,\y) -- node[right,font=\tiny\itshape,text=farmgray]{\lbl} (0,\y-0.7);

\end{tikzpicture}
\caption{Architecture système en 5 couches — Smart Farm AI v3.0}
\label{fig:arch_couches}
\end{figure}

\begin{tcolorbox}[infobox, title={Zone capture d'écran — Architecture déployée}]
\begin{center}
\begin{tikzpicture}
  \draw[dashed, farmgray, rounded corners=4pt, line width=1pt]
    (0,0) rectangle (14,4);
  \node[font=\large\bfseries\color{farmgray}] at (7,2.2)
    {[CAPTURE D'ÉCRAN]};
  \node[font=\small\color{farmgray}, align=center] at (7,1.3)
    {Diagramme Docker Compose — services backend, frontend,\\
     PostgreSQL, ChromaDB, Ollama et Caddy en production};
\end{tikzpicture}
\end{center}
\textit{Description :} Vue \code{docker ps} montrant les conteneurs actifs sur Oracle
Cloud ARM64. Insérer ici une capture du terminal avec les 6 services en statut
\code{Up}.
\end{tcolorbox}

\subsection{Composants de l'architecture}

\begin{table}[H]\centering
\caption{Description des composants principaux}
\label{tab:composants}
\begin{tabular}{lp{4cm}p{6cm}}
\toprule
\textbf{Composant} & \textbf{Technologie} & \textbf{Rôle} \\
\midrule
Owner Dashboard   & React 18 + Vite 5       & SPA principale : KPIs, télémétrie, apiculture, analytics \\
Worker PWA        & React + Dexie.js        & App mobile offline : scanner YOLO, tâches \\
FastAPI Backend   & Python 3.11 + FastAPI   & 50+ endpoints REST, WebSocket, JWT middleware \\
PostgreSQL        & PostgreSQL 15           & Base de données relationnelle (30+ tables) \\
ChromaDB          & ChromaDB 0.4+           & Base vectorielle pour RAG (cosine HNSW) \\
Ollama + Labess-7B & Ollama 0.1.x            & LLM local Darija tunisien \\
YOLO Models       & Ultralytics 8.1         & 8 modèles .pt en mémoire (lazy-cached) \\
IsolationForest   & scikit-learn 1.4        & Détection anomalies IoT (n=150, contam.=0.05) \\
DVC               & DVC 3.x                 & Versionnement données + pipeline reproductible \\
MLflow            & MLflow 2.x              & Tracking expériences, registre modèles \\
NODE\_A (ESP32)   & PlatformIO + C++        & Capteurs irrigation (sol, pression, débit) \\
NODE\_B (ESP32)   & PlatformIO + C++        & Capteurs rucher (poids, temp., humidité) \\
Docker Compose    & Docker 24               & Orchestration multi-services en production \\
Oracle Cloud      & ARM64 Free Tier         & Hébergement gratuit (LITE\_MODE sans GPU) \\
\bottomrule
\end{tabular}
\end{table}

% ─────────────────────────────────────────────────────────────
\section{Les besoins technologiques}
% ─────────────────────────────────────────────────────────────

\begin{table}[H]\centering
\caption{Besoins technologiques — Justification des choix}
\label{tab:tech_needs}
\small
\begin{tabular}{lp{3.5cm}p{4cm}l}
\toprule
\textbf{Besoin} & \textbf{Technologie retenue} & \textbf{Justification} & \textbf{Licence} \\
\midrule
API REST scalable    & FastAPI + Python 3.11   & Async natif, validation Pydantic, 50k req/s & MIT \\
Base relationnelle   & PostgreSQL 15           & ACID, multi-tenant, JSON natif & Open \\
Vision IA temps réel & YOLO v11 (Ultralytics)  & Meilleur compromis latence/mAP sur CPU & AGPL-3 \\
LLM dialecte local   & Labess-7B via Ollama    & Seul LLM Darija natif opérationnel & MIT \\
Base vectorielle     & ChromaDB (cosine HNSW)  & Simple à déployer, API Python, local & Apache 2 \\
Anomalie IoT         & IsolationForest (sklearn)& $O(n \log n)$, pas GPU requis & BSD \\
Versionnement ML     & DVC + MLflow            & Git-like pour données + tracking métriques & Apache 2 \\
Frontend SPA         & React 18 + Vite         & Écosystème riche, PWA, Recharts & MIT \\
Firmware IoT         & PlatformIO (C++)        & Compatible ESP32, multi-OS & Apache 2 \\
Simulation IoT       & Wokwi                   & Simulation ESP32 en ligne, HTTPS & Freemium \\
Déploiement          & Docker Compose + Oracle & Gratuit, ARM64, Caddy ACME auto & Gratuit \\
\bottomrule
\end{tabular}
\end{table}

% ─────────────────────────────────────────────────────────────
\section{Planification du projet}
% ─────────────────────────────────────────────────────────────

La planification du projet \farmAI{} v3.0 repose sur 6 sprints Scrum de 2 à 3
semaines chacun, couvrant 14 semaines au total. Chaque sprint cible un ensemble
cohérent de fonctionnalités avec un jalon de validation clair.

\begin{table}[H]\centering
\caption{Backlog de planification des 6 sprints}
\label{tab:planning}
\small
\begin{tabular}{ccp{2.2cm}p{5.5cm}p{2.3cm}}
\toprule
\textbf{Spr.} & \textbf{Durée} & \textbf{Focus} & \textbf{Livrables} & \textbf{Jalon} \\
\midrule
S1 & S0--S2 & Infra \& Auth & FastAPI (30+ routes), JWT+OTP, PostgreSQL 30+ tables, React scaffold & Auth OK \\
S2 & S2--S4 & IoT \& Télémétrie & NODE\_A/B firmware, \api{/iot/telemetry}, dashboard Recharts, Wokwi & IoT Live \\
S3 & S4--S7 & Apiculture & HiveIntel 7 sous-modules, CRUD, QR Codes, télémétrie NODE\_B & HiveIntel v1.0 \\
S4 & S7--S10 & Vision IA & 8 YOLO v11 intégrés, Scanner PWA (13 modèles), Map SIG Leaflet & YOLO Prod \\
S5 & S10--S12 & Assistant & Labess-7B RAG Darija, STT/TTS, LLaVA, Lite Mode Oracle & RAG Darija \\
S6 & S12--S14 & MLOps \& Final & DVC+MLflow, IsolationForest, Docker Compose, Oracle Cloud & Soutenance \\
\bottomrule
\end{tabular}
\end{table}

\begin{figure}[H]\centering
\begin{tikzpicture}[x=0.87cm, y=1cm]
\tikzset{
  sprint/.style={fill=#1, draw=#1!60!black, rounded corners=2pt, opacity=0.87},
  ms/.style={diamond, draw=farmgreen!80!black, fill=farmgreen!30,
    minimum size=0.5cm, inner sep=0pt},
  sout/.style={diamond, draw=red!70!black, fill=red!25,
    minimum size=0.6cm, inner sep=0pt}
}
\draw[->, thick, gray](-0.3,0)--(15.2,0) node[right, font=\small]{Semaines};
\foreach \w in {0,...,14}{
  \draw[gray!40,thin](\w,0.05)--(\w,-0.15);
  \node[below, font=\tiny\color{gray}] at (\w,-0.18){S\w};
}
\filldraw[sprint=farmblue!50]  (0,0.15)  rectangle (2,0.75);
\node[font=\tiny\bfseries\color{white}] at (1,0.45)   {S1 Infra};
\filldraw[sprint=farmblue!65]  (2,0.15)  rectangle (4,0.75);
\node[font=\tiny\bfseries\color{white}] at (3,0.45)   {S2 IoT};
\filldraw[sprint=farmblue!75]  (4,0.15)  rectangle (7,0.75);
\node[font=\tiny\bfseries\color{white}] at (5.5,0.45) {S3 Apiculture};
\filldraw[sprint=farmblue!85]  (7,0.15)  rectangle (10,0.75);
\node[font=\tiny\bfseries\color{white}] at (8.5,0.45) {S4 Vision IA};
\filldraw[sprint=farmblue!90]  (10,0.15) rectangle (12,0.75);
\node[font=\tiny\bfseries\color{white}] at (11,0.45)  {S5 Assistant};
\filldraw[sprint=farmblue!95]  (12,0.15) rectangle (14,0.75);
\node[font=\tiny\bfseries\color{white}] at (13,0.45)  {S6 MLOps};
\foreach \x/\lbl in {2/{Auth\\JWT}, 4/{IoT\\Live}, 7/{HiveIntel\\v1.0},
                      10/{YOLO\\Prod}, 12/{RAG\\Darija}}{
  \node[ms](m) at (\x,-0.7) {};
  \node[below=0.03cm of m, font=\tiny, align=center]{\lbl};
  \draw[dashed, gray!50, thin](\x,0)--(\x,-0.6);
}
\node[sout](ms) at (14,-0.7) {};
\node[below=0.03cm of ms, font=\tiny\bfseries, text=red!70!black, align=center]
  {Souten.\\v3.0};
\draw[dashed, red!50, thin](14,0)--(14,-0.6);
\end{tikzpicture}
\caption{Diagramme Gantt — 6 Sprints Scrum sur 14 semaines}
\label{fig:gantt_ch2}
\end{figure}

% ─────────────────────────────────────────────────────────────
\section{Conclusion}
% ─────────────────────────────────────────────────────────────

Ce chapitre a établi les fondements analytiques du projet. L'identification des
24 besoins fonctionnels et 12 besoins non-fonctionnels permet de définir un périmètre
précis et mesurable. La revue des travaux académiques confirme que \farmAI{} v3.0
répond à un manque réel dans la littérature agritech, en combinant vision par ordinateur,
IoT et LLM en Darija dans un seul système souverain. L'architecture en 5 couches et
la planification en 6 sprints constituent le cadre structurant du développement
présenté dans les chapitres suivants.
